\section{Platforms: documentation and local executions}\label{sec:platforms}

\subsection{S3 version identity}
An immutable version's identity and the history of which version is current
at a key are different objects. AWS documents two restorations: copying an old
version creates a new current version ID, while permanently deleting the
current version exposes its predecessor under the predecessor's original
ID~\cite{s3restore}. The current-key history $A/v_1\to B/v_2\to A/v_1$ is
therefore permitted although no ID was allocated twice. Deleting a delete marker
likewise restores the object below it, and $A/v_1\to\absent\to A/v_1$ remains a
counterexample when only successful reads are certifiable; the absent state
stays in the contract. A read that selects \code{versionId=v1} reads a
historical version, not the current key. These are documented operations, not
server faults, and neither the distinction between a version selector and a
condition on the current target~\cite{opendal} nor the ABA phenomenon is
claimed as new.

\subsection{Identity probes on nine implementations}
\code{experiments/reactivation_probes/} ran one fixed probe sequence on eight
S3-compatible implementations and a GCS emulator (Table~\ref{tab:probe-versions}):
\nProbeRequests{} loopback requests, every byte recorded, on Linux x86\_64 on
2026-11-05. Each server used a fresh temporary store and public fixture
credentials. The MinIO binary matches the release pinned for the earlier macOS
executions, so it also replicates them across platforms.

\begin{table}[t]
\caption{Implementations probed, their releases and the requests sent to each.}
\label{tab:probe-versions}
\small
\begin{tabular}{@{}llr@{}}
\toprule
Implementation & Version & Requests \\
\midrule
MinIO & \texttt{RELEASE.2025-09-07T16-13-09Z} & 51 \\
Ceph RADOS Gateway & \texttt{19.2.0} & 51 \\
SeaweedFS & \texttt{3.97} & 53 \\
Versity S3 Gateway & \texttt{1.10.0} & 51 \\
RustFS & \texttt{1.0.4} & 51 \\
Scality CloudServer & \texttt{9.0.27} & 51 \\
LocalStack & \texttt{4.9.2} & 52 \\
moto & \texttt{5.1.15} & 51 \\
fake-gcs-server & \texttt{1.56.1} & 34 \\
\bottomrule
\end{tabular}

\end{table}

All eight S3 implementations made the displaced version ID current again after
deleting the newest version by ID, and after inserting and removing a delete
marker (with a 404 in between). Over a COMPLIANCE-retained version, seven did
the same and refused direct deletion of the retained version (400 or 403);
SeaweedFS answered the retained PUT with 500, so that probe did not run. All
eight kept the identifier (\code{null} or absent) across a write of $A,B,A$
with versioning suspended. Copying the old version onto the key gave a fresh ID on seven, and the read
after the copy returned $A$'s first ETag again; LocalStack required replaced
metadata, and SeaweedFS refused all three request forms. Uploading the old bytes again gave a fresh ID on all eight
and repeated $A$'s ETag, and a conditional GET with $A$'s first ETag returned
304. No conditional GET was sent after a copy. The GCS emulator promoted
nothing when the live generation was deleted (404), and copy and upload
allocated new generations; this agrees with the GCS documentation but is not a
GCS measurement. \code{derive.py} reclassifies every cell from the raw records,
and Table~1 of the paper is generated in \code{table-artifacts/01-identity-probes/}
from a full copy of the traces.

\subsection{Two MinIO executions}
\paragraph{Figure~1a.} \code{experiments/versioned_objects/data/trace.json}
records six cases and 51 signed requests against an unmodified MinIO
RELEASE.2025-09-07T16-13-09Z (official binary, SHA-256 checked) on macOS; no
cloud account was used. Delete-current gives $A/\alpha, B/\beta, A/\alpha$;
copy restoration and re-upload give $A/\alpha, B/\beta, A/\gamma$; inserting and
removing a delete marker gives $A/\alpha,\absent,A/\alpha$ (Greek letters are
per-case aliases; the opaque IDs stay in the mapping file). A pinned control
keeps reading historical $A$ while the current key returns $B$. The sixth case
changes two keys in the order $A,B,B,A$, reading both stable keys after each
completed single-key mutation: the target pair (eligible, 10) is absent in all
five stable phases, yet the captures $A,B,B,A$ return matching IDs within each
key and induce overlapping bridges. This is the execution of Figure~1a.
\code{derive.py} verifies request and response hashes, selectors, the five
phase pairs and the capture order; the dedicated store and single writer are
execution assumptions that a transcript cannot prove.

\paragraph{Removing deletions.} \code{experiments/publication_profile/} records
eight cases, \nProfileRequests{} signed requests and five management commands
against the same server and a pinned management client. A non-owner publisher
may read current and historical objects, write, and delete without a version
ID, but is denied \code{DeleteObjectVersion} and every configuration change.

\begin{table*}[t]
\caption{The second MinIO execution: whether removing deletions removes reactivation.}
\label{tab:profile}
\small
\begin{tabularx}{\textwidth}{@{}lXX@{}}
\toprule
Case & Outcome & What it shows \\
\midrule
Owner under a bucket deny & DELETE of $B$ returns 204; $A$'s original ID returns & The owner's path is not constrained by that deny \\
Publisher deletes a data version & 403; $B$ stays current & Data-version deletion is blocked \\
Publisher deletes a marker & 403; the current read stays 404 & Marker removal is blocked \\
Publisher copies old $A$ & 200 and a fresh current ID & Restoration remains possible, with a fresh ID \\
Batch version deletion & 200 with two per-object AccessDenied errors & HTTP success alone does not classify a batch \\
Publisher changes configuration & 403 for versioning, policy and lifecycle & These changes are blocked \\
Two keys read in a stable phase & Equal IDs; bridges overlap for 793{,}750\,ns & A positive bundle under the declared contract \\
Multipart uploads completed in reverse & $B$ then $A$ current, both fresh IDs & One observed order, not AWS's \\
COMPLIANCE retention & Data deletion refused (400, WORM); marker insertion and removal 204 & Retained data still permits $A/\mathit{ID}\to\absent\to A/\mathit{ID}$ \\
\bottomrule
\end{tabularx}
\end{table*}

The rows of Table~\ref{tab:profile} are observations within one execution, not
independent trials. A first attempt stopped because the harness expected 403
for the retained-object deletion while the server returned 400 with a WORM
message; that raw attempt is preserved separately, and the reported run is a
fresh execution after the assertion was corrected. Privileged mutation, other
writers, lifecycle rules, replication, suspension of versioning and
restoration of the bucket or its authority are outside the profile, and
preventing every version deletion also prevents reclaiming noncurrent
versions.

\subsection{GCS generations}
For finalized objects at a fixed key in a continuously existing flat-namespace
bucket, generations are unique across successive data versions, although their
values need not increase; metadata-only updates keep the data generation and
change only the metageneration~\cite{gcsmetadata}. Deleting the live version
leaves no live version rather than promoting an older one, copying a
noncurrent version back allocates a new generation, and restoring a
soft-deleted object writes a new copy~\cite{gcsversioning}. Once a generation
stops being live at the key, these operations cannot make it live again, so
equal generations at two faithful current-name reads imply continuous stored
bytes in between, over any horizon in which generations are not reused. The
claim needs successful full current-name reads without a generation selector,
with the generation taken from the same response as the bytes; point fidelity
(no cache between client and origin, and a fixed encoding, since
\code{Accept-Encoding: gzip} is needed to prevent decompressive transcoding); a
fixed bucket, project and key; and the exclusion of bucket or project
restoration, hierarchical-namespace folder renames and appendable objects. It
says nothing about authorization, availability or response metadata, and no
GCS service was measured.

The official Python client pins versions silently. In
\code{experiments/gcs_selector/}, \code{google-cloud-storage} 3.17.1 runs
against a loopback fixture whose current data change from generation 101 (body
$A$) to 202 (body $B$). Reusing one \code{Blob} object for two reads sends
\code{generation=101} on the second and returns the old version; a fresh
\code{Blob} per read returns 202 and $B$; an explicit selector control behaves
as specified. This agrees with the client source~\cite{gcssdk}: the client
stores the response's generation and adds it as a selector to later downloads
of the same object. Each current-name observation must therefore use a fresh,
unbound object.

\subsection{Documented identifier semantics}
\code{research/platform-identifiers-round9.md} and its ledger record
\nAuditIdentifiers{} identifiers on \nAuditPlatforms{} storage, hosting, wiki,
package and HTTP-server platforms, with 257 quotations from \nAuditSources{}
official sources checked verbatim against copies retrieved on 2026-11-05. Each
entry records what a read or management interface exposes and what specific
restore, rollback or undelete operations do to it.
\nAuditReactivating{} identifiers are unique but can be made current again (S3
version IDs; B2 file IDs; DigitalOcean Spaces; Ceph RGW, from its tests; Swift
object versioning; SeaweedFS, from source; GCS hierarchical-namespace rename,
inferred; deploy and version IDs of Netlify, Vercel, Cloudflare Pages and
Firebase; npm dist-tags; MediaWiki revision IDs under deletion and
undeletion). \nAuditDerived{} derive from content or file attributes (MD5 and
SHA-1 ETags, content digests, Git and OCI identifiers, modification-time ETags
of nginx, Apache and Express), \nAuditNotUnique{} are not unique (S3
\code{null}, GCS metageneration) and \nAuditOther{} are unclear or
unversioned. The remaining \nAuditFresh{} receive fresh values on the restore
paths inspected (GCS flat-namespace generations, Azure Blob version IDs, Azure
ETags and snapshot promotion, Firebase release IDs, WordPress revision
restore). A fresh value on one path does not make a field sound for every
current read, and three corrections bound these entries.

\emph{MediaWiki.} Revision IDs are preserved across deletion and
undeletion~\cite{mediawikirevision}, so a page deleted and restored in full
returns as $A/r\to\absent\to A/r$; selective restoration can also make an older
revision current after another. Only ordinary edits and reverts, which
allocate new IDs, are inside the positive scope.

\emph{WordPress.} Restoring a revision creates a new one, but the current
post's \code{predecessor-version} link names the latest \emph{remaining}
revision row~\cite{wordpressrest}. For authorized edit-context reads and the
observation (\code{content.raw}, \code{predecessor-version.id}), the sequence
$A/r_1\to B/r_2\to A/r_3$, followed by deleting $r_3$ and $r_2$, returns to
$A/r_1$. This is derived from the official core source, not measured, and it
does not claim that the whole JSON response repeats.

\emph{Azure snapshots.} A snapshot's DateTime names an immutable historical
snapshot and is supplied later as a \code{snapshot} selector; the base-blob URI
and a version-selected URI are different reads~\cite{azuresnapshot}. Copying a
snapshot over the base blob writes a fresh current version and ETag, but the
snapshot DateTime is not a publication epoch on an unqualified current read.
Azure Blob version IDs are inside the positive scope only for versioned block
blobs; page and append blobs, whose updates need not create versions, and
point-in-time restore are outside it~\cite{azureversioning}.

Several classifications rest on project source or tests where the prose is
silent (Ceph, Swift, SeaweedFS, WordPress, MediaWiki undeletion, nginx), and
two on inference (GCS rename, Cloudflare Pages); the paper names only
explicitly documented cases. Documentation states intended behavior;
Table~1 of the paper tests implementations.
